\section*{Resumen}

La adhesión fibra--matriz es un factor determinante en el desempeño mecánico de los materiales
compuestos, ya que gobierna la transferencia de esfuerzos entre el refuerzo y la matriz. En este
contexto, el presente trabajo tuvo como objetivo implementar el ensayo Fiber-Bundle Pull-Out
(FBPO) como técnica de caracterización interfacial a mesoescala, mediante el diseño y fabricación
de un sistema experimental reproducible, con el fin de determinar la resistencia al corte
interfacial (IFSS aparente) y evaluar la utilidad del método como herramienta confiable para el
estudio de la adhesión fibra--matriz.

Para ello, se diseñó y fabricó un sistema experimental que incluyó moldes multicavidad, probetas
y dispositivos de sujeción, utilizando modelado CAD e impresión 3D, asegurando el control de
parámetros geométricos críticos como la longitud embebida y la alineación axial. El ensayo FBPO
fue implementado en probetas de material compuesto fabricadas con dos matrices epóxicas y dos
tipos de fibra de carbono, correspondientes a un sistema de grado aeroespacial y a una fibra
genérica. Se ejecutaron un total de 100 ensayos válidos bajo condiciones controladas, registrándose
las curvas fuerza--desplazamiento y determinándose la carga de despegue interfacial. La IFSS
aparente se estimó mediante un modelo de reducción basado en la relación entre la carga de
despegue y el área embebida.

Los resultados experimentales evidenciaron diferencias claras en los valores de IFSS aparente
entre los distintos sistemas fibra--matriz, tanto en magnitud como en dispersión. El análisis
estadístico, realizado mediante ANOVA de Welch y comparaciones múltiples de Games--Howell,
mostró que el tipo de fibra constituye el factor dominante en la respuesta interfacial dentro del
rango experimental estudiado, mientras que el efecto del tipo de matriz resultó secundario cuando
se empleó fibra genérica. La comparación con rangos de IFSS reportados en la literatura para
ensayos FBPO mostró una concordancia razonable, considerando las diferencias de escala,
materiales y condiciones experimentales.

En conjunto, los resultados confirman que el ensayo FBPO, acompañado de un diseño experimental
adecuado y un análisis estadístico robusto, constituye una técnica válida y reproducible para la
evaluación comparativa de la adhesión fibra--matriz a mesoescala, aportando al estudio y
comprensión del comportamiento interfacial en materiales compuestos.



